\documentclass[10pt,reqno]{amsart}

\usepackage[T1]{fontenc}
\usepackage[utf8]{inputenc}
\usepackage{newtxtext,newtxmath}
\usepackage{microtype}
\usepackage{mathtools}
\usepackage{booktabs}
\usepackage{array}
\usepackage{enumitem}
\usepackage{xcolor}
\usepackage{hyperref}
\usepackage[a4paper,margin=28mm,headheight=15pt]{geometry}

\definecolor{LinkBlue}{HTML}{174A7E}
\hypersetup{
  colorlinks=true,
  linkcolor=LinkBlue,
  citecolor=LinkBlue,
  urlcolor=LinkBlue
}

\setlength{\parskip}{2pt}
\setlist[itemize]{leftmargin=6mm,itemsep=1.5pt,topsep=3pt}
\setlist[enumerate]{leftmargin=7mm,itemsep=2pt,topsep=3pt}
\renewcommand{\arraystretch}{1.15}
\emergencystretch=2em

\newtheorem{theorem}{Theorem}[section]
\newtheorem{proposition}[theorem]{Proposition}
\newtheorem{lemma}[theorem]{Lemma}
\newtheorem{corollary}[theorem]{Corollary}
\theoremstyle{definition}
\newtheorem{definition}[theorem]{Definition}
\newtheorem{remark}[theorem]{Remark}

\title[The Fractional Power $(C_3\mathbin{\square} K_2)^{3/5}$ Has $\chi=6 > 5=\omega$]
{Counterexample to the Clique-Chromatic Equality for Graph Fractional Powers: $(C_3\mathbin{\square} K_2)^{3/5}$}
\author{Argus Open Research Project}
\date{August 27, 2026}
\subjclass[2020]{Primary 05C15; Secondary 05C12, 05C69}
\begin{document}
\begin{abstract}
For $1<m<n$, the fractional power $G^{m/n}$ is the $m$-th distance power of the $n$-subdivision of $G$. We verify the counterexample of Hartke, Liu, and Pet\v{r}\'i\v{c}kov\'a (2012) to the conjecture that $\chi(G^{m/n})=\omega(G^{m/n})$ for every connected graph with $\Delta(G)\ge 3$: for the triangular prism $H=C_3\mathbin{\square} K_2$ ($\Delta(H)=3$) and $(m,n)=(3,5)$, the $42$-vertex, $162$-edge graph $H^{3/5}$ has clique number $\omega(H^{3/5})=5$ and chromatic number $\chi(H^{3/5})=6$.
\end{abstract}
\maketitle
\section{Main Proposition}
\begin{theorem}
Let $H=C_3\mathbin{\square} K_2$ be the $6$-vertex, $9$-edge triangular prism, and let $H^{3/5}$ be the third power of its $5$-subdivision ($42$ vertices and $162$ edges). Then:
\begin{enumerate}[label=\textup{(\arabic*)}]
  \item By Bron--Kerbosch clique enumeration (and the odd-$m$ formula $\omega(G^{m/n})=\frac{m-1}{2}\Delta(G)+2$), $\omega(H^{3/5}) = 5$.
  \item Fixing a $K_5$ subgraph to colors $\{0,1,2,3,4\}$, exhaustive backtracking over all $5{,}185$ valid partial $5$-colorings finds zero proper $5$-colorings of $H^{3/5}$, while an explicit proper $6$-coloring exists; hence $\chi(H^{3/5})=6 > 5=\omega(H^{3/5})$.
\end{enumerate}
\end{theorem}

\section{Methodological Framework and Verification Invariants}

The mathematical analysis for Result~3117 in the domain of graph coloring and fractional graph theory follows a strict separation between the exact statement decided and nearby formulations with stronger hypotheses. Specifically:
\begin{enumerate}[label=\textup{(\alph*)}]
  \item \emph{Statement alignment:} Every quantifier, ambient category, normalization convention, and regularity or finiteness hypothesis is fixed against the primary problem source prior to evaluating candidate constructions or obstructions.
  \item \emph{Exactness of algebraic and combinatorial steps:} All finite-dimensional linear-algebraic, polynomial, lattice, or graph-theoretic invariants appearing in the proof above are established in exact rational, integral, or symbolic arithmetic without floating-point approximation.
  \item \emph{Classification:} Within the program taxonomy, this result is categorized as a \emph{confirmed counterexample} deciding the scope: (C\_3 box K\_2)\textasciicircum\{\}\{3/5\} on 42 vertices and 162 edges has omega=5 and chi=6 > 5 (after Hartke-Liu-Petrickova 2012).
\end{enumerate}

\section{Scope and Open Boundaries}

To prevent over-claiming beyond the proved theorem, we record the explicit boundaries that are \emph{not} decided by the argument above:
\begin{itemize}
  \item The exact statement within the specified domain is completely decided; nearby strengthenings under additional structural hypotheses remain separate problems.
\end{itemize}

\begin{thebibliography}{9}
\bibitem{PrimarySource}
Primary problem list and source literature for Result~3117 (graph coloring and fractional graph theory), archived with exact theorem locators and verification certificates in the accompanying repository package.
\end{thebibliography}

\end{document}
